\section*{Hoofdstuk \ref{ch:b1:nucleophilic-substitution} --- Nucleofiele substitutie}

\begin{solution}{exo:b1:nucleophilic-substitution:1}
\ce{CH3CH2OH + Br-}; \ce{CH3OCH2CH3 + I-}; \ce{CH3CH2CH2CN + Cl-};
\ce{CH3NH3+ + Br-}. \omterm{def:b1:nucleophilic-substitution:substitution}{Substraten}: de halogeenalkanen; \omterm{def:b1:electronic-effects:nucleophile}{nucleofielen}:
\ce{OH-}, \ce{CH3CH2O-}, \ce{CN-}, \ce{NH3}; \omterm{def:b1:electronic-effects:nucleophile}{vertrekkende groepen}: de
halogenide-ionen.
\end{solution}

\begin{solution}{exo:b1:nucleophilic-substitution:2}
Tweede orde: verdrievoudigd; beide gehalveerd: gedeeld door 4. Eerste
orde in alleen RY: ongewijzigd; gehalveerd.
\end{solution}

\begin{solution}{exo:b1:nucleophilic-substitution:3}
SN2 (primair, goed \omterm{def:b1:electronic-effects:nucleophile}{nucleofiel}, aprotisch); SN1 (tertiair, water); SN2
(methyl, sterk \omterm{def:b1:electronic-effects:nucleophile}{nucleofiel}); SN1-solvolyse (secundair, zwak neutraal
\omterm{def:b1:electronic-effects:nucleophile}{nucleofiel}, \omterm{def:b1:intermolecular-forces-solvents:solvent-classes}{protisch oplosmiddel}), mogelijk met wat SN2.
\end{solution}

\begin{solution}{exo:b1:nucleophilic-substitution:4}
Het cyanide neemt de plaats tegenover het broom in en de drie andere
groepen klappen om. In het product staat de CN-groep, zoals eerst Br,
op de eerste plaats, en de andere rangen blijven ongewijzigd: de
omgekeerde schikking heeft de tegengestelde descriptor,
\cip{S}-2-methylbutaannitril.
\end{solution}

\begin{solution}{exo:b1:nucleophilic-substitution:5}
1. \omterm{def:b1:electronic-effects:cleavage}{Heterolyse}: pijl van de C--Br-binding naar Br, wat \ce{(CH3)3C+} en
\ce{Br-} geeft (traag). 2. Een \omterm{def:b1:lewis-resonance-vsepr:lewis}{vrij paar} van water naar het kationische
koolstofatoom: \ce{(CH3)3C-OH2+}. 3. Een \omterm{def:b1:lewis-resonance-vsepr:lewis}{vrij paar} van een ander
watermolecuul neemt een proton van de \ce{OH2+}-groep, waarbij het
O--H-paar naar de zuurstof teruggaat. Alleen de eerste, trage stap
bepaalt de snelheid; water, het oplosmiddel, is in zo'n overmaat dat zijn
concentratie niet verandert.
\end{solution}

\begin{solution}{exo:b1:nucleophilic-substitution:6}
Het koolstofatoom dat Br draagt, is primair, maar ernaast vult een
tert-butylgroep de ruimte achter de C--Br-binding: de \omterm{def:b1:elementary-steps:energy-profile}{overgangstoestand}
van SN2 is erg vol. SN1 zou een primair \omterm{def:b1:electronic-effects:intermediates}{carbokation} vereisen, dat te
instabiel is.
\end{solution}

\begin{solution}{exo:b1:nucleophilic-substitution:7}
Jodide vervangt jodide volgens SN2, telkens met inversie: \cip{S} wordt
\cip{R} tot beide in gelijke mate aanwezig zijn, een \omterm{def:b1:stereochemistry-in-depth:enantiomers}{racemisch mengsel}.
Elke substitutie zet één \cip{S}-molecuul om in één \cip{R}: de rotatie
daalt met de bijdrage van twee moleculen, dus de rotatie neemt twee keer
zo snel af als het \omterm{def:b1:nucleophilic-substitution:substitution}{substraat} gesubstitueerd wordt.
\end{solution}

\begin{solution}{exo:b1:nucleophilic-substitution:8}
SN2: broommethaan $>$ 1-broompropaan $>$ 2-broompropaan $>$
2-broom-2-methylpropaan (ruimtelijke hindering). SN1: de omgekeerde
volgorde (stabiliteit van het \omterm{def:b1:electronic-effects:intermediates}{carbokation}), waarbij broommethaan en
1-broompropaan nauwelijks reageren.
\end{solution}

\begin{solution}{exo:b1:nucleophilic-substitution:9}
$x_{\text{inv}} - x_{\text{ret}} = 0.12$ en $x_{\text{inv}} + x_{\text{ret}} = 1$:
\qty{56}{\percent} geïnverteerd, \qty{44}{\percent} met behoud van
\omterm{def:b1:stereochemistry-in-depth:isomers}{configuratie}. Vooral SN1-racemisatie, met een lichte overmaat aan
inversie: het vertrekkende bromide schermt zijn kant nog af wanneer het
water aankomt.
\end{solution}

\begin{solution}{exo:b1:nucleophilic-substitution:10}
$v = k_1k_2[\mathrm{RY}][\mathrm{Nu}]/(k_{-1}[\mathrm{Y^-}] + k_2[\mathrm{Nu}])$
(\cref{prop:b1:nucleophilic-substitution:sn1-law}). \ce{Y-} toevoegen
vergroot de noemer: meer \omterm{def:b1:electronic-effects:intermediates}{carbokationen} keren terug naar het \omterm{def:b1:nucleophilic-substitution:substitution}{substraat} in
plaats van ingevangen te worden, en de snelheid daalt. Verwaarloosbaar als
$k_2[\mathrm{Nu}] \gg k_{-1}[\mathrm{Y^-}]$, zoals wanneer het \omterm{def:b1:electronic-effects:nucleophile}{nucleofiel}
het oplosmiddel is.
\end{solution}

\begin{solution}{exo:b1:nucleophilic-substitution:11}
Gelijke concentraties: $-\mathrm{d}c/\mathrm{d}t = kc^2$, $1/c = 1/C_0 + kt$,
$t_{1/2} = 1/(kC_0) = \qty{5.0e5}{s}$ (ongeveer zes dagen). Met een
vijftigvoudige overmaat hydroxide is $[\ce{OH-}] \approx 50C_0$ constant:
$c = C_0
e^{-k't}$ met $k' = 50kC_0 = \qty{1.0e-4}{s^{-1}}$, $t_{1/2} = \ln 2/k' =
\qty{6.9e3}{s}$, ongeveer twee uur.
\end{solution}

\begin{solution}{exo:b1:nucleophilic-substitution:12}
Bromide zou \ce{OH-} moeten verdrijven, een \omterm{def:b1:predominant-reaction:strong-acid}{sterke base} en een slechte
\omterm{def:b1:electronic-effects:nucleophile}{vertrekkende groep}: geen reactie. In geconcentreerd HBr wordt de OH
geprotoneerd tot \ce{-OH2+}, waarvan de \omterm{def:b1:electronic-effects:nucleophile}{vertrekkende groep} water is.
Tertiaire alcohol: water vertrekt, wat het \omterm{def:b1:electronic-effects:intermediates}{carbokation} geeft, dat door
\ce{Br-} ingevangen wordt (SN1). Primaire alcohol: \ce{Br-} valt het
koolstofatoom van \ce{R-OH2+} van achteren aan terwijl water vertrekt
(SN2).
\end{solution}

\begin{solution}{pb:b1:nucleophilic-substitution:1}
\textbf{1.} Elk molecuul dat reageert, geeft één \ce{H+} en één \ce{Cl-};
het geleidingsvermogen is de som van de ionaire bijdragen, evenredig met
de hoeveelheid ionen.
\textbf{2.} $[\mathrm{RCl}]/[\mathrm{RCl}]_0 = (\kappa_\infty -
\kappa)/\kappa_\infty$.
\textbf{3.} $\ln[\mathrm{RCl}] = \ln[\mathrm{RCl}]_0 - kt$: zet
$\ln(\kappa_\infty - \kappa)$ uit tegen $t$.
\textbf{4.} 0.693, 0.603, 0.513, 0.423, 0.333, 0.153, $-0.117$, $-0.387$.
\textbf{5.} Ja, overal een richtingscoëfficiënt van $-0.0180$ per minuut:
eerste orde.
\textbf{6.} Elke orde in water: de concentratie ervan is constant
(ontaarding van de orde).
\textbf{7.} $k = 0.0180/60 = \qty{3.0e-4}{s^{-1}}$.
\textbf{8.} $t_{1/2} = \ln 2/k = \qty{2.3e3}{s}$, \qty{38.5}{min}.
\textbf{9.} $\ln 10/k = \qty{7.7e3}{s}$, \qty{128}{min}.
\textbf{10.} $\ln(\kappa_\infty - \kappa)$: 0.693 bij $t = 0$, 0.109 bij
\qty{10}{min}, enzovoort: richtingscoëfficiënt $-0.0584$ per minuut,
$k = \qty{9.7e-4}{s^{-1}}$.
\textbf{11.} $9.7/3.0 = 3.2$.
\textbf{12.} Bij \qty{30}{min}: $2.000(1 - e^{-0.0584 \times 30}) = 1.653$,
gemeten 1.654.
\textbf{13.} SN1: tertiair \omterm{def:b1:nucleophilic-substitution:substitution}{substraat}, \omterm{def:b1:intermolecular-forces-solvents:solvent-classes}{protisch oplosmiddel}, neutraal
\omterm{def:b1:electronic-effects:nucleophile}{nucleofiel}.
\textbf{14.} Zoals in \cref{exo:b1:nucleophilic-substitution:5}, met Cl.
\textbf{15.} Bij de trage ionisatie is alleen het \omterm{def:b1:nucleophilic-substitution:substitution}{substraat} betrokken:
$v =
k[\mathrm{RCl}]$.
\textbf{16.} Hydroxide zou alleen in de snelle invangstap kunnen
optreden, na de snelheidsbepalende ionisatie; de snelheid hangt er niet
van af.
\textbf{17.} Het \omterm{def:b1:electronic-effects:intermediates}{carbokation} is vlak en \omterm{def:b1:stereochemistry-in-depth:stereocentre}{achiraal}; water addeert even
gemakkelijk aan beide zijden: racemische alcohol.
\textbf{18.} Een hoge eerste barrière (ionisatie) naar het \omterm{def:b1:electronic-effects:intermediates}{carbokation} in
een ondiep dal, daarna een lage barrière naar het invangen: de eerste
stap is snelheidsbepalend.
\textbf{19.} Het zou een primair \omterm{def:b1:electronic-effects:intermediates}{carbokation} vormen (te instabiel), en
water is een te zwak \omterm{def:b1:electronic-effects:nucleophile}{nucleofiel} voor een snelle SN2.
\textbf{20.} $k = A\,e^{-E_a/RT}$.
\textbf{21.} $E_a = R\ln(k_2/k_1)/(1/T_1 - 1/T_2)$.
\textbf{22.} $E_a = 8.314 \times \ln(9.74/3.00)/(1/298.15 - 1/308.15) =
\qty{9.0e4}{J/mol}$.
\textbf{23.} De hoogte van de barrière van de ionisatie, de
\omterm{def:b1:elementary-steps:rds}{snelheidsbepalende stap}.
\textbf{24.} $k = 3.0 \times 10^{-4} \exp\bigl(\frac{90\,000}{8.314}
(\frac{1}{298.15} - \frac{1}{318.15})\bigr) = \qty{2.9e-3}{s^{-1}}$.
\textbf{25.} \textbf{$E_a = \qty{90}{kJ/mol}$}.
\end{solution}
